% Parte 1: fenómeno, marco teórico, métodos y metodología.
% Se incluye desde informe.tex; no compilar este archivo solo.

% ===========================================================================
\section{Explicación del fenómeno}
% ===========================================================================

Un cuerpo que cae en el aire está sometido a dos fuerzas: su peso, que es constante, y la resistencia del aire (arrastre), que se opone al movimiento y crece con la rapidez. Al inicio de la caída el cuerpo se mueve despacio, el arrastre es pequeño y el cuerpo acelera casi como lo haría en el vacío. A medida que gana rapidez el arrastre aumenta, hasta que llega a igualar al peso. En ese momento la fuerza neta es cero, la aceleración desaparece y el cuerpo sigue cayendo a rapidez constante. Esa rapidez límite se llama \emph{velocidad terminal} y se denota $\vt$.

Este comportamiento tiene consecuencias prácticas importantes. Un paracaidista en caída libre no acelera indefinidamente, sino que se estabiliza alrededor de los \qty{180}{km/h}; una gota de lluvia llega al suelo a pocos metros por segundo en lugar de hacerlo a la rapidez que tendría tras caer cientos de metros en el vacío. También explica por qué el resultado del experimento atribuido a Galileo, en el que objetos de distinta masa caen al mismo tiempo, solo se cumple de forma aproximada en presencia de aire.

En este proyecto se estudia esta dinámica en una dimensión (solo movimiento vertical). Se comparan dos leyes de arrastre, una lineal y una cuadrática, junto con la caída en el vacío como referencia, en dos situaciones:
\begin{enumerate}
  \item un paracaidista en caída libre desde \qty{4000}{m};
  \item cuatro pelotas de distintos deportes soltadas desde \qty{56}{m}.
\end{enumerate}
Además, se valida la implementación numérica con un análisis de convergencia de los dos métodos de integración utilizados.

% ===========================================================================
\section{Marco teórico}
% ===========================================================================

\subsection{Ecuación de movimiento}

Se toma el eje $y$ positivo hacia arriba y se describe el estado del cuerpo con su altura $y$ y su velocidad $v$. Por la segunda ley de Newton, el movimiento queda descrito por el sistema de dos ecuaciones diferenciales ordinarias de primer orden
\begin{equation}
  \frac{dy}{dt} = v, \qquad \frac{dv}{dt} = -g + \frac{F_d(v)}{m},
  \label{eq:movimiento}
\end{equation}
donde $m$ es la masa del cuerpo, $g$ la aceleración de la gravedad y $F_d(v)$ la fuerza de arrastre. Como el cuerpo cae, $v<0$ durante todo el movimiento.

\subsection{Modelos de arrastre}

Se consideran tres modelos para $F_d$, resumidos en el Cuadro~\ref{tab:modelos}. En cada caso, la velocidad terminal se obtiene al igualar el arrastre con el peso, es decir, al imponer $\dot v = 0$.

\begin{table}[H]
  \centering
  \caption{Modelos de arrastre considerados.}
  \label{tab:modelos}
  \renewcommand{\arraystretch}{1.6}
  \begin{tabular}{@{}llll@{}}
    \toprule
    Modelo & Fuerza $F_d(v)$ & Ecuación para $v$ & Velocidad terminal \\
    \midrule
    Vacío & $0$ & $\dot v = -g$ & no existe \\
    Lineal (Stokes) & $-b\,v$ & $\dot v = -g - \dfrac{b}{m}\,v$ & $\vt = \dfrac{mg}{b}$ \\
    Cuadrático & $-k\,v\lvert v\rvert$ & $\dot v = -g - \dfrac{k}{m}\,v\lvert v\rvert$ & $\vt = \sqrt{\dfrac{mg}{k}}$ \\
    \bottomrule
  \end{tabular}
\end{table}

En el modelo cuadrático, la constante de arrastre es
\begin{equation}
  k = \tfrac12 \rho\, C_d\, A,
  \qquad\text{de modo que}\qquad
  \vt = \sqrt{\frac{2mg}{\rho\, C_d\, A}},
  \label{eq:k}
\end{equation}
donde $\rho$ es la densidad del aire, $C_d$ el coeficiente de arrastre y $A$ el área frontal del cuerpo. Se usa $v\lvert v\rvert$ en lugar de $v^2$ para que la fuerza siempre se oponga al movimiento, sin importar el signo de $v$.

\subsection{Validez de cada modelo: número de Reynolds}

El modelo adecuado depende del régimen del flujo alrededor del cuerpo, que se caracteriza con el número de Reynolds
\begin{equation}
  Re = \frac{\rho\, \lvert v\rvert\, L}{\mu},
  \label{eq:reynolds}
\end{equation}
donde $L$ es una longitud característica (el diámetro, en el caso de una esfera) y $\mu$ la viscosidad dinámica del aire. El arrastre lineal solo es válido en el régimen de Stokes, $Re \lesssim 1$, que corresponde a partículas microscópicas como polvo o gotas de niebla. Para objetos macroscópicos en aire, en el rango $10^3 < Re < 2\times10^5$, el coeficiente $C_d$ de una esfera es aproximadamente constante (cercano a 0.47) y el arrastre es cuadrático~\cite{white}.

Para dimensionar esto, la rapidez a la que $Re = 1$ es $v = \mu/(\rho L)$. Para una pelota de ping-pong esa rapidez es de \qty{0.00037}{m/s} y para una bola de boliche de \qty{0.000068}{m/s}. Cualquier caída real de estos objetos ocurre a rapideces varios órdenes de magnitud mayores, así que el modelo lineal no les aplica. En este proyecto el modelo lineal se usa como punto de comparación y para validar los métodos numéricos, ya que tiene una solución exacta sencilla.

\subsection{Soluciones analíticas}
\label{sec:analiticas}

Los tres modelos tienen solución cerrada, lo que permite medir con exactitud el error de los métodos numéricos. Se denota $h_0$ la altura inicial y $v_0 \le 0$ la velocidad inicial; en todos los escenarios de este trabajo $v_0 = 0$.

\paragraph{Vacío.}
\begin{equation}
  v(t) = v_0 - g t, \qquad y(t) = h_0 + v_0 t - \tfrac12 g t^2 .
\end{equation}
Desde el reposo, el tiempo de caída es $t = \sqrt{2h_0/g}$ y la rapidez de impacto es $\sqrt{2gh_0}$.

\paragraph{Arrastre lineal.} Con $\tau = m/b$,
\begin{align}
  v(t) &= -\vt + (v_0 + \vt)\, e^{-t/\tau}, \\
  y(t) &= h_0 - \vt t + \tau\,(v_0 + \vt)\left(1 - e^{-t/\tau}\right).
\end{align}
La rapidez se acerca a $\vt$ de forma exponencial. Desde el reposo, el tiempo para alcanzar el \qty{95}{\percent} de $\vt$ es $t_{95} = \tau \ln 20$.

\paragraph{Arrastre cuadrático.} Para $-\vt < v_0 \le 0$, con $\tau_0 = \atanh(\lvert v_0\rvert/\vt)$~\cite{taylor},
\begin{align}
  v(t) &= -\vt \tanh\!\left(\frac{g t}{\vt} + \tau_0\right), \\
  y(t) &= h_0 - \frac{\vt^2}{g}\,
          \ln\frac{\cosh\!\left(\frac{g t}{\vt} + \tau_0\right)}{\cosh \tau_0}.
\end{align}
Desde el reposo, el tiempo para alcanzar el \qty{95}{\percent} de $\vt$ y el tiempo de caída hasta el suelo son
\begin{equation}
  t_{95} = \frac{\vt}{g}\,\atanh(0.95),
  \qquad
  t_{\text{imp}} = \frac{\vt}{g}\,\acosh\!\left(e^{g h_0/\vt^2}\right).
\end{equation}

En caídas largas, $\cosh x$ excede el rango de los números de punto flotante. Por eso, en la implementación se evalúa $\ln\cosh x$ mediante la identidad equivalente y numéricamente estable
\begin{equation}
  \ln\cosh x = \lvert x\rvert + \ln\!\left(1 + e^{-2\lvert x\rvert}\right) - \ln 2 .
\end{equation}

% ===========================================================================
\section{Métodos de simulación}
% ===========================================================================

El sistema~\eqref{eq:movimiento} se escribe en forma vectorial con el estado $\bs = (y, v)$ y $\bff(\bs) = (v,\ a(v))$, donde $a(v) = -g + F_d(v)/m$. Para integrarlo se implementaron dos métodos de paso fijo, sin recurrir a bibliotecas como \texttt{scipy}~\cite{burden}.

\subsection{Euler explícito}

Es el método más simple: avanza el estado usando la derivada evaluada al inicio del paso,
\begin{equation}
  \bs_{n+1} = \bs_n + \Delta t\, \bff(\bs_n).
\end{equation}
Su error global es $O(\Delta t)$, es decir, es de orden~1.

\subsection{Runge-Kutta de orden 4}

RK4 combina cuatro evaluaciones de la derivada dentro de cada paso:
\begin{gather}
  \bk_1 = \bff(\bs_n), \quad
  \bk_2 = \bff\!\left(\bs_n + \tfrac{\Delta t}{2}\bk_1\right), \quad
  \bk_3 = \bff\!\left(\bs_n + \tfrac{\Delta t}{2}\bk_2\right), \quad
  \bk_4 = \bff(\bs_n + \Delta t\,\bk_3), \\
  \bs_{n+1} = \bs_n + \frac{\Delta t}{6}\left(\bk_1 + 2\bk_2 + 2\bk_3 + \bk_4\right).
\end{gather}
Su error global es $O(\Delta t^4)$. En la práctica, al dividir el paso entre 10 el error de Euler disminuye 10 veces y el de RK4 disminuye \num{10000} veces. A cambio, RK4 requiere 4 evaluaciones de $\bff$ por paso en lugar de una; medido en nuestra implementación, cada paso de RK4 tarda entre 4.5 y 5 veces más que uno de Euler.

\subsection{Detección del impacto con el suelo}

La integración avanza hasta el primer paso en el que $y \le 0$. Como el impacto ocurre, en general, entre dos pasos, su instante y su velocidad se estiman por interpolación lineal entre el último punto con $y > 0$ y el primero con $y \le 0$:
\begin{equation}
  s = \frac{y_n}{y_n - y_{n+1}}, \qquad
  t_{\text{imp}} = t_n + s\,\Delta t, \qquad
  v_{\text{imp}} = v_n + s\,(v_{n+1} - v_n).
\end{equation}

\subsection{Implementación}

El simulador está escrito en Python, usando únicamente \texttt{numpy} para los arreglos y \texttt{matplotlib} para las figuras. El código se organiza en el paquete \texttt{freefall}:
\begin{itemize}
  \item \texttt{models.py}: cuerpos (masa, $C_d$, área, longitud característica) y modelos de arrastre;
  \item \texttt{analytic.py}: soluciones exactas de la Sección~\ref{sec:analiticas};
  \item \texttt{solvers.py}: pasos de Euler y RK4 e integración con detección del impacto;
  \item \texttt{scenarios.py}: parámetros de los escenarios y cálculo de métricas;
  \item \texttt{convergence.py}: errores contra la solución analítica y orden observado;
  \item \texttt{plots.py}, \texttt{report.py} y \texttt{\_\_main\_\_.py}: figuras, tablas e interfaz de línea de comandos.
\end{itemize}
Cada modelo de arrastre es una clase que implementa la aceleración $a(v)$ y la velocidad terminal, de modo que los integradores son independientes del modelo. El Listado~\ref{lst:solvers} muestra el núcleo de los integradores.

\begin{lstlisting}[caption={Pasos de Euler y RK4 y detección del impacto (\texttt{solvers.py}).}, label={lst:solvers}]
def euler_step(f, state, dt):
    return state + dt * f(state)

def rk4_step(f, state, dt):
    k1 = f(state)
    k2 = f(state + 0.5 * dt * k1)
    k3 = f(state + 0.5 * dt * k2)
    k4 = f(state + dt * k3)
    return state + dt / 6 * (k1 + 2 * k2 + 2 * k3 + k4)

# Dentro del ciclo de integración:
new = step(f, states[-1], dt)
if stop_at_ground and new[0] <= 0:
    old = states[-1]
    s = old[0] / (old[0] - new[0])
    impact_time = times[-1] + s * dt
    impact_velocity = old[1] + s * (new[1] - old[1])
\end{lstlisting}

% ===========================================================================
\section{Metodología}
% ===========================================================================

\subsection{Parámetros generales}

El Cuadro~\ref{tab:constantes} resume las constantes físicas y numéricas comunes a todos los escenarios. Salvo que se indique lo contrario, todas las simulaciones usan RK4 con $\Delta t = \qty{0.01}{s}$ y parten del reposo.

\begin{table}[H]
  \centering
  \caption{Constantes usadas en todos los escenarios.}
  \label{tab:constantes}
  \begin{tabular}{@{}ll@{}}
    \toprule
    Constante & Valor \\
    \midrule
    Gravedad $g$ & \qty{9.81}{m/s^2} \\
    Densidad del aire $\rho$ & \qty{1.225}{kg/m^3} (nivel del mar, \qty{15}{\degreeCelsius}) \\
    Viscosidad dinámica del aire $\mu$ & \qty{1.81e-5}{\pascal\second} \\
    Paso de tiempo por defecto $\Delta t$ & \qty{0.01}{s} \\
    Método por defecto & RK4 \\
    \bottomrule
  \end{tabular}
\end{table}

\subsection{Escenario 1: paracaidista}

Se simula un salto desde \qty{4000}{m} en posición de vientre abajo, sin abrir el paracaídas, desde el reposo hasta el suelo. Los parámetros se muestran en el Cuadro~\ref{tab:paracaidista-param}. Se comparan los tres modelos: vacío, lineal y cuadrático.

\begin{table}[H]
  \centering
  \caption{Parámetros del paracaidista.}
  \label{tab:paracaidista-param}
  \begin{tabular}{@{}ll@{}}
    \toprule
    Parámetro & Valor \\
    \midrule
    Masa $m$ & \qty{80}{kg} \\
    Coeficiente de arrastre $C_d$ & 1.0 \\
    Área frontal $A$ & \qty{0.5}{m^2} \\
    Longitud característica $L$ (para $Re$) & \qty{1.8}{m} \\
    $k = \tfrac12\rho C_d A$ & \qty{0.30625}{kg/m} \\
    Velocidad terminal (cuadrático) & \qty{50.62}{m/s} (\qty{182}{km/h}) \\
    $b$ del modelo lineal & \qty{15.503}{kg/s} \\
    $\tau = m/b = \vt/g$ & \qty{5.160}{s} \\
    $Re$ a velocidad terminal & \num{6.2e6} \\
    \bottomrule
  \end{tabular}
\end{table}

Una decisión de diseño importante es la elección de $b$. Su valor no se tomó de la física del problema (el modelo lineal no aplica a un paracaidista, cuyo $Re$ es del orden de $10^6$), sino que se \emph{calibró para que la velocidad terminal fuera la misma que en el modelo cuadrático}: $b = mg/\vt$. De esta forma las curvas de ambos modelos terminan en el mismo valor y la comparación aísla únicamente la diferencia de forma entre una ley de arrastre lineal y una cuadrática. Con esta calibración, además, la escala de tiempo $m/b$ del modelo lineal coincide con $\vt/g$, la escala natural del modelo cuadrático (\qty{5.16}{s} en ambos casos).

\subsection{Escenario 2: comparación de pelotas}

Se sueltan desde el reposo cuatro pelotas reales a \qty{56}{m} de altura, la altura aproximada de la Torre de Pisa, en referencia al experimento atribuido a Galileo. Se usa el modelo cuadrático con $C_d = 0.47$ (esfera lisa) y la caída en el vacío como referencia. Las masas y los radios corresponden a los valores reglamentarios de cada deporte (la bola de boliche es de 16~lb). El Cuadro~\ref{tab:pelotas-param} muestra los parámetros y la velocidad terminal resultante.

\begin{table}[H]
  \centering
  \small
  \caption{Parámetros de las pelotas ($C_d = 0.47$).}
  \label{tab:pelotas-param}
  \begin{tabular}{@{}lrrrrrr@{}}
    \toprule
    Pelota & \thead{Masa\\(kg)} & \thead{Radio\\(m)} & \thead{Diámetro\\(m)} & \thead{Área $A$\\(m$^2$)} & \thead{$k$\\(kg/m)} & \thead{$\vt$\\(m/s)} \\
    \midrule
    Ping-pong & 0.0027 & 0.020  & 0.040  & 0.001257 & 0.000362 & 8.56 \\
    Tenis     & 0.058  & 0.0335 & 0.067  & 0.003526 & 0.001015 & 23.68 \\
    Béisbol   & 0.145  & 0.0366 & 0.0732 & 0.004208 & 0.001211 & 34.27 \\
    Boliche   & 7.26   & 0.108  & 0.216  & 0.036644 & 0.010549 & 82.17 \\
    \bottomrule
  \end{tabular}
\end{table}

La pregunta que busca responder este escenario es qué tanto importa el aire en una caída corta, dependiendo de la relación entre la masa y el área de cada pelota.

\subsection{Escenario 3: validación y convergencia}
\label{sec:diseno-convergencia}

Para validar los métodos y medir su orden de convergencia se diseñó el siguiente experimento:
\begin{itemize}
  \item \textbf{Caso de prueba:} el paracaidista del Escenario~1 con los modelos lineal y cuadrático, integrado con Euler y con RK4.
  \item \textbf{Intervalo:} se integra en $[0, 20]$~s \emph{sin detener la simulación en el suelo}. Así el error medido corresponde solo al integrador y no se mezcla con el de la interpolación del impacto.
  \item \textbf{Pasos de tiempo:} 13 valores de $\Delta t$: 1, 0.5, 0.25, 0.1, 0.05, 0.025, 0.01, 0.005, 0.0025, 0.001, 0.0005, 0.00025 y 0.0001~s.
  \item \textbf{Métrica:} para cada paso se mide el error absoluto máximo en $y$ y en $v$ contra la solución analítica.
  \item \textbf{Orden observado:} es la pendiente de un ajuste lineal de $\log(\text{error})$ contra $\log(\Delta t)$. Se descartan el paso de \qty{1}{s} y los puntos con error menor a $10^{-9}$, donde domina el error de redondeo.
  \item \textbf{Orden local:} también se calcula el orden entre pasos consecutivos,
  $p_i = \dfrac{\log(e_{i-1}/e_i)}{\log(\Delta t_{i-1}/\Delta t_i)}$.
\end{itemize}
Adicionalmente, se repiten los escenarios 1 y 2 con Euler para medir el efecto del método en los resultados, y se integra el modelo cuadrático con un paso grande ($\Delta t = \qty{2}{s}$) para ilustrar visualmente la diferencia entre ambos métodos.

